%%%%%%%%%%%%%%%%%%%%%%%%%%%%%%%%%%%%%%%%%%%%%%%%%%%%%%%%%%%%%%%%%%%%%%%
% How to use writeLaTeX: 
%
% You edit the source code here on the left, and the preview on the
% right shows you the result within a few seconds.
%
% Bookmark this page and share the URL with your co-authors. They can
% edit at the same time!
%
% You can upload figures, bibliographies, custom classes and
% styles using the files menu.
%
%%%%%%%%%%%%%%%%%%%%%%%%%%%%%%%%%%%%%%%%%%%%%%%%%%%%%%%%%%%%%%%%%%%%%%
\documentclass[a4paper]{article}

\usepackage{renote-template}
\usepackage{graphicx,url}
\usepackage[brazil]{babel}
\usepackage[utf8]{inputenc}  
\usepackage{longtable}
\usepackage[none]{hyphenat}
\usepackage{fancyhdr}
\usepackage{setspace}
\usepackage{hyperref}
\usepackage[alf, bibjustif, abnt-etal-list=5, abnt-emphasize=bf,
            abnt-etal-text=it, abnt-repeated-author-omit=no        ]{abntex2cite}

\fancyhf{}
\fancyhead[C]{\thepage}
\pagestyle{fancy}
\sloppy
\singlespacing

\title{Título em Português (Renote Jorunal)}

\author{
    Autor 1, Afiliação 1, E-Mail 1 \\ 
    Autor 2, Afiliação 2, E-Mail 2 \\
    Autor 3, Afiliação 3, E-Mail 3
}

\begin{document} 

\maketitle

\begin{resumo} 
O artigo devem possuir obrigatoriamente um resumo, no seu idioma original, com 100 a 200 palavras ou 10 linhas no máximo; seguido de um conjunto de palavras-chaves (de 3 até 5, escritas em minúsculas), também no idioma original do artigo. Logo a seguir deve aparecer a tradução do resumo e das palavras-chaves, segundo o seguinte critério: se o idioma do artigo for o português ou espanhol, a tradução será para o inglês; se o idioma do artigo for o inglês, a tradução será para o português. (Abstract, Keywords); se for espanhol (Resumen) ou inglês (Abstract), a tradução será para o português. O resumo deve ser claro, sucinto e explicar, em um único parágrafo, o(s) objetivo(s) pretendido(s), procurando justificar sua importância (sem incluir referências bibliográficas), os principais procedimentos adotados, os resultados mais expressivos e conclusões.
\end{resumo}

\begin{palavas-chave}
    Palavra Chave 1, Palavra Chave 2, Palavra Chave 3, Palavra Chave 4, Palavra Chave 5.
\end{palavas-chave}

\begin{englishtitle}
    English Title (Renote Journal)
\end{englishtitle}

\begin{abstract}
O artigo devem possuir obrigatoriamente um resumo, no seu idioma original, com 100 a 200 palavras ou 10 linhas no máximo; seguido de um conjunto de palavras-chaves (de 3 até 5, escritas em minúsculas), também no idioma original do artigo. Logo a seguir deve aparecer a tradução do resumo e das palavras-chaves, segundo o seguinte critério: se o idioma do artigo for o português ou espanhol, a tradução será para o inglês; se o idioma do artigo for o inglês, a tradução será para o português. (Abstract, Keywords); se for espanhol (Resumen) ou inglês (Abstract), a tradução será para o português. O resumo deve ser claro, sucinto e explicar, em um único parágrafo, o(s) objetivo(s) pretendido(s), procurando justificar sua importância (sem incluir referências bibliográficas), os principais procedimentos adotados, os resultados mais expressivos e conclusões.
\end{abstract}

\begin{keywords}
    Palavra Chave 1, Palavra Chave 2, Palavra Chave 3, Palavra Chave 4, Palavra Chave 5.
\end{keywords}

\section{Estrutura}

Tamanho da folha: para  impressão em papel  A4 (21,0 cm x 29,7 cm) sendo observadas as seguintes margens: superior 3,0; cm inferior 2,0 cm; esquerda e direita 3,0 cm. Alinhar a margem direita evitando separações silábicas com barras ou outros sinais.

Utilizar fonte "Times" tamanho 12 ou equivalente. Não utilizar fontes que ocupam muito espaço tal como "Bookman".  No texto utilizar espaço 1.  O fim de uma seção e o cabeçalho da próxima são separados por um espaço extra.

Todas as páginas do trabalho devem ser numeradas, com exceção da folha de rosto.  Os números, em algarismos arábicos, são colocados, no centro da margem superior. Sugere-se utilizar o comando "cabeçalho" do editor de texto para especificar a paginação. Entre a paginação e o texto deixar algum espaço para destacar os dois elementos. 

\section{Teste}
O número de tabelas e figuras utilizadas no artigo deve ser limitado a compreensão e elucidação do texto. Devem ser inseridas no corpo do texto, para identificação da sua posição e do tamanho aproximado (Figura~\ref{fig:exampleFig1}).

\subsection{Subseção}
O número de tabelas e figuras utilizadas no artigo deve ser limitado a compreensão e elucidação do texto. Devem ser inseridas no corpo do texto, para identificação da sua posição e do tamanho aproximado (Figura~\ref{fig:exampleFig1}).

\section{Figures, Tables and Captions}\label{sec:figs}

O número de tabelas e figuras utilizadas no artigo deve ser limitado a compreensão e elucidação do texto. Devem ser inseridas no corpo do texto, para identificação da sua posição e do tamanho aproximado (Figura~\ref{fig:exampleFig1}).

Tabelas e figuras possuem numeração independente, que deve ser feita seqüencialmente na ordem em que são citadas no texto. Devem também ter uma legenda auto-explicativa, sendo que as tabelas terão legendas na parte superior e as figuras as terão na parte inferior, centralizadas em relação à tabela ou figura. A legenda inicia com o termo "Tabela" ou "Figura" (primeira letra em maiúscula), de acordo com o caso, seguido de um espaço e do número de ordem seqüencial, em algarismos arábicos, seguido de hífen entre espaços e do texto da legenda, com a primeira letra da primeira palavra em maiúscula e as demais em minúscula, exceção feitas àquelas que normalmente são escritas em maiúsculas. Devem ser citadas no texto como "Tabela \ref{tab:exTable1}" e "Figura" seguidas de espaço e do número correspondente.

\begin{figure}[ht]
\centering
\includegraphics[width=.5\textwidth]{fig1.jpg}
\caption{A typical figure}
\label{fig:exampleFig1}
\end{figure}

\begin{figure}[ht]
\centering
\includegraphics[width=.3\textwidth]{fig2.jpg}
\caption{This figure is an example of a figure caption taking more than one
  line and justified considering margins mentioned in Section~\ref{sec:figs}.}
\label{fig:exampleFig2}
\end{figure}

\begin{table}[ht]
\centering
\caption{Variables to be considered on the evaluation of interaction
  techniques}
\label{tab:exTable1}
\includegraphics[width=.7\textwidth]{table.jpg}
\end{table}

\section{Agradecimentos}
Usados apenas quando necessário. Devem ser sucintos, diretos e dirigidos apenas a pessoas ou instituições que contribuíram substancialmente para a elaboração do trabalho.

\section{Referências}

As referências a autores e fontes são inseridas no texto colocando entre parênteses o sobrenome do autor, com inicial maiúscula, seguida da data, conforme o exemplo: \cite{smith:99}. Quando for conhecida a paginação, pode-se incluí-la, por exemplo, (Tarouco, 1991, p. 237). Havendo mais de um título do(s) autor(es) no mesmo ano, deve-se distingui-las utilizando uma letra minúscula (a,b,c) depois da data, por exemplo, (Tarouco, 1991b). Quando houver três ou mais autores, no texto será citado apenas o primeiro autor seguido de "et al.", mas nas referências bibliográficas, no final do artigo, os demais nomes também deverão aparecer. Quando o nome do autor é citado diretamente no texto pode-se colocar entre parênteses apenas a data de publicação, por exemplo,

\citeonline{boulic:91}. Na citação de citação, identifica-se a obra diretamente consultada; o autor e/ou a obra citada nesta é assim indicado: \cite{knuth:84} citada por Vicari \cite{teste}.

As referências completas de cada autor e fonte citadas no texto devem aparecer no final do artigo sob o título "Referências Bibliográficas", ordenadas alfabeticamente pelo(s) sobrenomes(s) do(s) autore(s).

Os autores do artigo devem certificar-se de que as referências citadas no texto contam da lista de referências com datas exatas e nomes de autores corretamente grafados. A exatidão dessas referências é de responsabilidade dos autores do artigo. Comunicações pessoais, trabalhos inéditos ou em andamento poderão ser citados quando absolutamente necessários, mas não devem ser incluídos na lista de referências bibliográficas; apenas citados no texto ou nas Notas do texto.

\addto{\captionsbrazil}{% caso use \usepackage[brazil]{babel}
\renewcommand{\bibname}{Refer\^{e}ncias Teses}
}
\bibliographystyle{renote}
\bibliography{References/references}

\end{document}
